% PDF page 292
\source{gilbarg-trudinger:page-0292}

\begin{corollary}
\label{gt-cor-11-2}
\dcref{Corollary 11.2}
\group{gt-ch11}
\source{gilbarg-trudinger:page-0292}
Corollary 11.2. \emph{Let $\G$ be a closed convex set in a Banach space $\B$ and let $T$ be a
continuous mapping of $\G$ into itself such that the image $T\G$ is precompact. Then T has
a fixed point.}
\end{corollary}

\begin{theorem}
\label{gt-thm-11-3}
\dcref{Theorem 11.3}
\group{gt-ch11}
\source{gilbarg-trudinger:page-0292}
\uses{gt-cor-11-2}
Theorem 11.3. \emph{Let $T\!$ be a compact mapping of a Banach space $\B$ into itself, and
suppose there exists a constant $M$ such that}

\[(11.2)\qquad \|x\|_{\B} < M\]

\emph{for all $x \in \B$ and $\sigma \in [0,1]$ satisfying $x = \sigma Tx$. Then $T$ has a fixed point.}

Proof. We can assume without loss of generality that $M = 1$. Let us define a
mapping $T^*$ by
\end{theorem}
\begin{proof}
\uses{gt-cor-11-2}
\[(11.1)\qquad \|J_k x - x\| \le \frac{\sum \dist(x,\G-B^i)\|x_i-x\|}{\sum \dist(x,\G-B^i)} < \frac{1}{k}.\]

The mapping $J_k \circ T$ when restricted to $\G_k$ is accordingly a continuous mapping of
$\G_k$ into itself and hence, by virtue of the Brouwer fixed point theorem, possesses a
fixed point $x^{(k)}$. (Note that $\G_k$ is homeomorphic to a closed ball in some Euclidean
space.) Since $\G$ is compact, a subsequence of the sequence $x^{(k)}$, $k = 1, 2, \ldots$, con-
verges to some $x \in \G$. We claim that x is a fixed point of $T$. For, applying (11.1) to
$Tx^{(k)}$, we have

\[
\|x^{(k)} - Tx^{(k)}\| = \|J_k \circ T x^{(k)} - Tx^{(k)}\| < \frac{1}{k},
\]

and, since $T$ is continuous, we infer $Tx = x$. $\square$

As will be demonstrated in the next chapter, Theorem 11.1 is applicable to
broad classes of equations in two variables. For later purposes we note the following
extension of Theorem 11.1.

Corollary 11.2. \emph{Let $\G$ be a closed convex set in a Banach space $\B$ and let $T$ be a
continuous mapping of $\G$ into itself such that the image $T\G$ is precompact. Then T has
a fixed point.}

In the above theorems we note an essential difference from the contraction
mapping principle, Theorem 5.1, in that the fixed points whose existence is asserted
are not necessarily unique.

11.2. The Leray-Schauder Theorem: a Special Case

A continuous mapping between two Banach spaces is called \emph{compact} (or \emph{completely}
\emph{continuous}) if the images of bounded sets are precompact (that is, their closures are
compact). The following consequence of Corollary 11.2 is the fixed point result
most often applied in our approach to the Dirichlet problem for quasilinear
equations.

Theorem 11.3. \emph{Let $T\!$ be a compact mapping of a Banach space $\B$ into itself, and
suppose there exists a constant $M$ such that}

\[(11.2)\qquad \|x\|_{\B} < M\]

\emph{for all $x \in \B$ and $\sigma \in [0,1]$ satisfying $x = \sigma Tx$. Then $T$ has a fixed point.}

Proof. We can assume without loss of generality that $M = 1$. Let us define a
mapping $T^*$ by

% PDF page 293
\source{gilbarg-trudinger:page-0293}


\noindent\textbf{11.2. The Leray-Schauder Theorem: a Special Case}

\[
T^{*}x=
\begin{cases}
Tx & \text{if }\|Tx\|\leq 1,\\
\dfrac{Tx}{\|Tx\|} & \text{if }\|Tx\|\geq 1.
\end{cases}
\]

\noindent Then $T^{*}$ is a continuous mapping of the closed unit ball $\bar B$ in $\B$ into itself. Since $T\bar B$
is precompact the same is true of $T^{*}\bar B$. Hence by Corollary 11.2 the mapping $T^{*}$
has a fixed point $x$. We claim that $x$ is also a fixed point of $T$. For, suppose that
$\|Tx\|\geq 1$. Then $x=T^{*}x=\sigma Tx$ if $\sigma=1/\|Tx\|$, and $\|x\|=\|T^{*}x\|=1$, which contra-
dicts (11.2) with $M=1$. Hence $\|Tx\|<1$ and consequently $x=T^{*}x=Tx$. $\square$

\medskip
\noindent\textit{Remark.}\quad Theorem 11.3 implies that if $T$ is any compact mapping of a Banach
space into itself (whether or not (11.2) holds), then for some $\sigma\in(0,1]$ the mapping
$\sigma T$ possesses a fixed point. Furthermore, if the estimate (11.2) holds then $\sigma T$ has a
fixed point for all $\sigma\in[0,1]$.
\end{proof}

\medskip
\noindent In order to apply Theorem 11.3 to the Dirichlet problem for quasilinear
equations, we fix a number $\beta\in(0,1)$ and take the Banach space $\B$ to be the Hölder
space $C^{1,\beta}(\bar\Omega)$, where $\Omega$ is a bounded domain in $\R^n$. Let $\Q$ be the operator given by

\begin{equation}
\tag{11.3}
\Q u=a^{ij}(x,u,\D u)\D_{ij}u+b(x,u,\D u)
\end{equation}

\noindent and assume that $\Q$ is elliptic in $\bar\Omega$, that is, the coefficient matrix $[a^{ij}(x,z,p)]$ is positive for all
$(x,z,p)\in\bar\Omega\times\R\times\R^n$. We also assume, for some $\alpha\in(0,1)$, that the
coefficients $a^{ij},\,b\in C^{1}(\bar\Omega\times\R\times\R^n)$, that the boundary $\partial\Omega\in C^{2,\alpha}$ and that $\varphi$ is a
given function in $C^{2,\alpha}(\bar\Omega)$. For all $v\in C^{1,\beta}(\bar\Omega)$, the operator $T$ is defined by letting
$u=Tv$ be the unique solution in $C^{2,\alpha\beta}(\bar\Omega)$ of the linear Dirichlet problem,

\begin{equation}
\tag{11.4}
a^{ij}(x,v,\D v)\D_{ij}u+b(x,v,\D v)=0\ \text{in }\Omega,\qquad u=\varphi\ \text{on }\partial\Omega.
\end{equation}

\noindent The unique solvability of the problem (11.4) is guaranteed by the linear existence
result, Theorem 6.14. The solvability of the Dirichlet problem, $\Q u=0$ in $\Omega$, $u=\varphi$
on $\partial\Omega$, in the space $C^{2,\alpha}(\bar\Omega)$ is thus equivalent to the solvability of the equation
$u=Tu$ in the Banach space $\B=C^{1,\beta}(\bar\Omega)$. The equation $u=\sigma Tu$ in $\B$ is equivalent
to the Dirichlet problem

\begin{equation}
\tag{11.5}
\Q_{\sigma}u=a^{ij}(x,u,\D u)\D_{ij}u+\sigma b(x,u,\D u)=0\ \text{in }\Omega,\qquad u=\sigma\varphi\ \text{on }\partial\Omega.
\end{equation}

\noindent By applying Theorem 11.3, we can then prove the following criterion for existence.

\medskip
\begin{theorem}
\label{gt-thm-11-4}
\dcref{Theorem 11.4}
\group{gt-ch11}
\source{gilbarg-trudinger:page-0293,gilbarg-trudinger:page-0294}
\uses{gt-thm-11-3}
\textbf{Theorem 11.4.}\quad \textit{Let $\Omega$ be a bounded domain in $\R^n$ and suppose that $\Q$ is elliptic in}
\newline
\textit{$\bar\Omega$ with coefficients $a^{ij},\,b\in C^{1}(\bar\Omega\times\R\times\R^n)$, $0<\alpha<1$. Let $\partial\Omega\in C^{2,\alpha}$ and $\varphi\in C^{2,\alpha}(\bar\Omega)$.}
\newline
\textit{Then, if for some $\beta>0$ there exists a constant $M$, independent of $u$ and $\sigma$, such that}
\newline
\textit{every $C^{2,\alpha}(\bar\Omega)$ solution of the Dirichlet problems, $\Q_{\sigma}u=0$ in $\Omega$, $u=\sigma\varphi$ on $\partial\Omega$,}
\newline
\textit{$0\leq\sigma\leq 1$, satisfies}

\begin{equation}
\tag{11.6}
\|u\|_{C^{1,\beta}(\bar\Omega)}<M,
\end{equation}
\end{theorem}

% PDF page 294
\source{gilbarg-trudinger:page-0294}
\begin{gather*}
\textit{it follows that the Dirichlet problem, } Qu=0 \text{ in } \Omega,\ u=\varphi \text{ on } \partial\Omega,\ \text{is solvable in } C^{2,\alpha}(\overline{\Omega}).
\end{gather*}

\begin{proof}
\uses{gt-thm-11-4}
\textit{Proof.} In view of the remarks preceding the statement of the theorem, it only remains to show that the operator $T$ is continuous and compact. By virtue of the global Schauder estimate, Theorem 6.6, $T$ maps bounded sets in $C^{1,\beta}(\overline{\Omega})$ into bounded sets in $C^{2,\alpha\beta}(\overline{\Omega})$ which (by Arzela's theorem) are precompact in $C^2(\overline{\Omega})$ and $C^{1,\beta}(\overline{\Omega})$. In order to show the continuity of $T$, we let $v_m$, $m=1,2,\ldots$, converge to $v$ in $C^{1,\beta}(\overline{\Omega})$. Then, since the sequence $\{Tv_m\}$ is precompact in $C^2(\overline{\Omega})$, every subsequence in turn has a convergent subsequence. Let $\{Tv_m\}$ be such a convergent subsequence with limit $u\in C^2(\overline{\Omega})$. Then since

\begin{gather*}
a^{ij}(x,v,Dv)D_{ij}u+b(x,v,Dv) \\
= \lim_{m\to\infty}\{a^{ij}(x,\bar v_m,D\bar v_m)D_{ij}T\bar v_m+b(x,\bar v_m,D\bar v_m)\}=0,
\end{gather*}

\noindent we must have $u=Tv$, and hence the sequence $\{Tv_m\}$ itself converges to $u$. \qed
\end{proof}

\section*{11.3. \ An Application}

Theorem 11.4 reduces the solvability of the Dirichlet problem $Qu=0$ in $\Omega$, $u=\varphi$ on $\partial\Omega$ to the apriori estimation in the space $C^{1,\beta}(\overline{\Omega})$, for some $\beta>0$, of the solutions of a related family of problems. In practice it is desirable to break the derivation of the apriori estimates into four stages:

\begin{enumerate}
\item Estimation of $\sup_\Omega |u|$;
\item Estimation of $\sup_{\partial\Omega} |Du|$ in terms of $\sup_\Omega |u|$;
\item Estimation of $\sup_\Omega |Du|$ in terms of $\sup_{\partial\Omega} |Du|$ and $\sup_\Omega |u|$;
\item Estimation of $[Du]_{\beta;\Omega}$, for some $\beta>0$, in terms of $\sup_\Omega |Du|$, $\sup_\Omega |u|$.
\end{enumerate}

Step I has already been treated in Chapter 10; (see Theorems 10.3, 10.4 and 10.9). Steps II and III are treated in Chapters 14 and 15. In Chapter 13 it will be shown that Step IV can be carried out under very general hypotheses on $Q$. We shall illustrate the overall procedure here by considering a problem where the required estimates are readily obtained from some of the results in earlier chapters. Namely, let us suppose that either $Q$ has the special divergence form,

\begin{equation}
(11.7)\qquad Qu=\operatorname{div} A(Du),
\end{equation}

\noindent or that $n=2$ and $Q$ has the form

\begin{equation}
(11.8)\qquad Qu=a^{ij}(x,u,Du)D_{ij}u,\qquad i,j=1,2.
\end{equation}

% PDF page 295
\source{gilbarg-trudinger:page-0295}

11.3. An Application

We shall demonstrate in Chapter 14 that geometric conditions on the boundary $\partial\Omega$
play an important role in the solvability of the Dirichlet problem for quasilinear
equations. For our present purposes, we shall require that the boundary manifold

\[
\Gamma=(\partial\Omega,\varphi)=\{(x,z)\in \partial\Omega\times\mathbb{R}\mid z=\varphi(x)\}
\]

satisfies a bounded slope condition, that is, for every point $P\in\Gamma$ there exist two
planes in $\mathbb{R}^{n+1}$, $z=\pi_P^-(x)$ and $z=\pi_P^+(x)$, passing through $P$ such that:

\begin{align*}
\text{(i)}\quad &\pi_P^-(x)\le \varphi(x)\le \pi_P^+(x) \qquad \forall x\in \partial\Omega;\\
\text{(ii)}\quad &\text{the slopes of these planes are uniformly bounded, independently of $P$,}\\
&\text{by a constant $K$; that is }|D\pi_P^\pm|\le K \quad \text{for all }P\in\Gamma.
\end{align*}

If $\partial\Omega\in C^2$, $\varphi\in C^2(\overline{\Omega})$ and $\partial\Omega$ is uniformly convex (that is, its principal curvatures
are bounded away from zero), then $\Gamma$ satisfies a bounded slope condition ; (see
[HA]). We now can assert the following existence result.

\textbf{Theorem 11.5.} \textit{Let $Q$ have either of the forms (11.7) or (11.8), and suppose that
$Q$, $\Omega$ and $\varphi$ satisfy the hypotheses of Theorem 11.4. Then if, in addition, the boundary
manifold $(\partial\Omega,\varphi)$ satisfies the bounded slope condition, it follows that the Dirichlet
problem, $Qu=0$ in $\Omega$, $u=\varphi$ on $\partial\Omega$, is solvable in $C^2(\overline{\Omega})$.}

\textit{Proof.} Since $Q_\sigma=Q$, we must estimate the solutions of the Dirichlet problems,
$Qu=0$ in $\Omega$, $u=\sigma\varphi$ on $\partial\Omega$, $0\le \sigma\le 1$. Let us take in turn the different stages described
above.

\begin{enumerate}
\item[I.] From the weak maximum principle, (Theorems 3.1 or 10.3), we have

\begin{equation*}
\tag{11.9}
\sup_\Omega |u|=\sigma\sup_{\partial\Omega}|\varphi|\le \sup_{\partial\Omega}|\varphi|.
\end{equation*}

\item[II.] The bounded slope condition provides linear barriers which serve to
estimate $Du$ on $\partial\Omega$. For clearly

\[
a^{ij}(x,u,Du)D_{ij}\pi_P^\pm=0,
\]

and hence, by the weak maximum principle,

\[
\sigma\pi_P^-(x)\le u(x)\le \sigma\pi_P^+(x),
\]

for all $x\in\Omega$. Consequently we have

\begin{equation*}
\tag{11.10}
\sup_{\partial\Omega}|Du|\le \sigma K\le K,
\end{equation*}

where $K$ is the assumed bound for the slopes of $\pi_P^\pm$.
\item[III.] Steps III and IV will follow from the fact that the derivatives $D_k u$, $k=
1,\dots,n$, are weak solutions of simple linear divergence structure equations of the
\end{enumerate}


% PDF page 296
\source{gilbarg-trudinger:page-0296}
type treated in Chapter 8. Let us first suppose that $Q$ has the form (11.7) and write
the equation $Qu=0$ in its integral form,

\begin{equation}
\int_{\Omega} A(\Du)\cdot D\eta \, dx = 0 \qquad \forall \eta \in C_0^1(\Omega).
\tag{11.11}
\end{equation}

Fixing $k$, replacing $\eta$ by $D_k\eta$ and then integrating by parts, we obtain

\begin{equation}
\int_{\Omega} D_j A^i(\Du) D_k u D_i\eta \, dx = 0 \qquad \forall \eta \in C_0^1(\Omega),
\end{equation}

or, if $\w = D_k u$,

\begin{equation}
\int_{\Omega} a^{ij}(\Du) D_j \w D_i\eta \, dx = 0 \qquad \forall \eta \in C_0^1(\Omega).
\end{equation}

The function $\w \in C^1(\overline{\Omega})$ is thus a weak solution of the linear elliptic equation

\begin{equation}
D_i\bigl(\bar a^{ij}(x)D_j \w\bigr)=0,
\tag{11.12}
\end{equation}

where $\bar a^{ij}(x)=a^{ij}(\Du(x))$, and hence, by the weak maximum principle in Section 3.6
(see also Theorem 8.1), we have

\begin{equation}
\sup_{\Omega} |Du| = \sup_{\partial\Omega} |Du| \le K.
\tag{11.13}
\end{equation}

Next, if $Q$ has the form (10.8), the equation $Qu=0$ is equivalent to

\begin{equation}
\frac{a^{11}}{a^{22}}D_{11}u + \frac{2a^{12}}{a^{22}}D_{12}u + D_{22}u = 0,
\end{equation}

so that

\begin{equation}
\int_{\Omega}\left(\frac{a^{11}}{a^{22}}D_{11}u + \frac{2a^{12}}{a^{22}}D_{12}u + D_{22}u\right)\eta \, dx = 0 \qquad \forall \eta \in C_0^1(\Omega).
\end{equation}

Replacing $\eta$ by $D_1\eta$ and integrating by parts, we obtain for $\w = D_1u$,

\begin{equation}
\int_{\Omega}\left\{\left(\frac{a^{11}}{a^{22}}D_1\w + \frac{2a^{12}}{a^{22}}D_2\w\right)D_1\eta + D_2\w D_2\eta\right\} dx = 0,
\end{equation}

and hence $\w$ is a weak solution of the linear elliptic equation

\begin{equation}
D_i(a^{ij}_1(x)D_j\w)=0, \qquad i,j=1,2.
\tag{11.14}
\end{equation}


% PDF page 297
\source{gilbarg-trudinger:page-0297}

\pagestyle{empty}


with coefficient matrix
\[
[a^{ij}_1(x)] =
\begin{bmatrix}
\dfrac{a^{11}}{u^{22}}(x, u(x), Du(x)) &
\dfrac{2a^{12}}{u^{22}}(x, u(x), Du(x)) \\
0 & 1
\end{bmatrix}.
\]

Similarly, it follows that $D_2u$ is a weak solution of a corresponding linear elliptic equation and hence again, by the weak maximum principle, the estimate (11.13) holds.

IV. The equations (11.12) and (11.14) for the derivatives $D_k u$ will satisfy the hypotheses of Theorem 8.24 with constants $\lambda$ and $\Lambda$ depending on $\sup_{\Omega}|u|$, $\sup_{\Omega}|Du|$ and the coefficients $a^{ij}$. Consequently we obtain an interior H\"older estimate for $Du$, that is, for any subdomain $\Omega' \Subset \Omega$, we have
\begin{equation}
[Du]_{\beta;\Omega'} \le C d^{-\beta}
\tag{11.15}
\end{equation}
where the positive constants $C$ and $\beta$ are independent of $u$ and $\sigma$, and $d=\operatorname{dist}(\Omega',\partial\Omega)$. We cannot however infer a global H\"older estimate for $Du$ directly from the results of Chapter 8. Instead we proceed as follows. We first use the smoothness of $\partial\Omega$ to map portions of $\partial\Omega$ into the hyperplane $x_n=0$. The derivatives $D_{y_k}u$, $k=1,\ldots,n-1$, with respect to the transformed coordinates $y_1,\ldots,y_n$ can then be estimated by means of Theorem 8.29. The remaining derivative $D_{y_n}u$ is finally estimated by using the equation itself together with Morrey's estimate, Theorem 7.19. The details of this procedure are carried out in Chapter 13 for the general divergence structure equation. The resulting estimate
\begin{equation}
[Du]_{\beta;\Omega} \le C
\tag{11.16}
\end{equation}
with positive constants $\beta$ and $C$ independent of $u$ and $\sigma$ completes the proof of Theorem 11.5. $\square$

We note here that from the results of Chapter 14 it will follow that the bounded slope condition in the hypothesis of Theorem 11.5 can be replaced by the boundedness of the quantity
\[
\frac{\Lambda(x,z,p)|p|}{\mathcal{G}(x,z,p)}
\]
for $x \in \overline{\Omega}$, $|z| \le \sup_{\partial\Omega} |\varphi|$, $|p| \ge 1$; (see Theorem 14.1). Furthermore, if $\Omega$ is convex, the bounded slope condition can be replaced by the boundedness of the quantity
\[
\frac{\Lambda(x,z,p)}{a^{ij}(x,z,p)(p_i - D_i\varphi)(p_j - D_j\varphi)}
\]
for $x \in \overline{\Omega}$, $|z| \le \sup_{\partial\Omega} |\varphi|$, $|p| \ge 1$ (see Theorem 14.2).


% PDF page 298
\source{gilbarg-trudinger:page-0298}



\noindent 11.4. \ The Leray-Schauder Fixed Point Theorem

\medskip

For certain applications it is desirable to replace the family of Dirichlet problems,
$Q_{\sigma}u=0$ in $\Omega$, $u=\sigma\varphi$ on $\partial\Omega$, $0\leq \sigma\leq 1$, used in Theorem 11.4 by other families
which depend differently on the parameter $\sigma$. Accordingly, we shall require the
following generalization of Theorem 11.3.

\medskip

\noindent \textbf{Theorem 11.6.} \ \emph{Let $\B$ be a Banach space and let $T$ be a compact mapping of $\B\times[0,1]$
into $\B$ such that $T(x,0)=0$ for all $x\in\B$. Suppose there exists a constant $M$ such that}

\[
\tag{11.17}
\|x\|_{\B}<M
\]

\emph{for all $(x,\sigma)\in\B\times[0,1]$ satisfying $x=T(x,\sigma)$. Then the mapping $T_{1}$ of $\B$ into itself
given by $T_{1}x=T(x,1)$ has a fixed point.}

\medskip

Theorem 11.6 will be derived from the following consequence of Corollary 11.2.

\medskip

\noindent \textbf{Lemma 11.7.} \ \emph{Let $B=B_{1}(0)$ denote the unit ball in $\B$ and let $T$ be a continuous
mapping of $\overline{B}$ into $\B$ such that $T\overline{B}$ is precompact and $T\partial B\subset B$. Then $T$ has a fixed
point.}

\medskip

\noindent \textit{Proof.} \ We define a mapping $T^{*}$ by

\[
T^{*}x=
\begin{cases}
Tx & \text{if }\|Tx\|\leq 1,\\[0.4em]
\dfrac{Tx}{\|Tx\|} & \text{if }\|Tx\|\geq 1.
\end{cases}
\]

Clearly $T^{*}$ is a continuous mapping of $\overline{B}$ into itself and since $T\overline{B}$ is precompact, the
same is true of $T^{*}\overline{B}$. Hence, by Corollary 11.2, $T^{*}$ has a fixed point $x$ and since
$T\partial B\subset B$ we must have $\|x\|<1$ and therefore $x=Tx$.

\medskip

\noindent \textit{Proof of Theorem 11.6.} \ We can assume without loss of generality that $M=1$.
For $0<\varepsilon\leq 1$, let us define a mapping $T^{*}$ from $\overline{B}$ into $\B$ by

\[
T^{*}x=T_{\varepsilon}^{*}x=
\begin{cases}
T\!\left(\dfrac{x}{\|x\|},\dfrac{1-\|x\|}{\varepsilon}\right)
& \text{if }1-\varepsilon\leq \|x\|\leq 1,\\[1.1em]
T\!\left(\dfrac{x}{1-\varepsilon},1\right)
& \text{if }\|x\|<1-\varepsilon.
\end{cases}
\]

The mapping $T^{*}$ is clearly continuous, $T^{*}\overline{B}$ is precompact by the compactness of $T$
and $T^{*}\partial B=0$. Hence, by Lemma 11.7, the mapping $T^{*}$ has a fixed point $x(\varepsilon)$.
We now set

\[
\varepsilon=\frac{1}{k},\qquad x_{k}=x\!\left(\frac{1}{k}\right),\qquad \sigma_{k}=
\begin{cases}
k(1-\|x_{k}\|) & \text{if }1-\dfrac{1}{k}\leq \|x_{k}\|\leq 1,\\[0.8em]
1 & \text{if }\|x_{k}\|<1-\dfrac{1}{k}.
\end{cases}
\]


% PDF page 299
\source{gilbarg-trudinger:page-0299}



\noindent where $k=1,2,\dots$. By the compactness of $T$ we can assume, by passing to a
subsequence if necessary, that the sequence $\{(x_k,\sigma_k)\}$ converges in $\B\times[0,1]$ to
$(x,\sigma)$. It then follows that $\sigma=1$. For if $\sigma<1$, we must have $\|x_k\|\geq 1-1/k$ for
sufficiently large $k$ and hence $\|x\|=1$, $x=T(x,\sigma)$, which contradicts (11.17).
Since $\sigma=1$, we then have by the continuity of $T$ that $\Tstar_{1/k}x_k\to T(x,1)$, and therefore
$x$ is also a fixed point of $T_1$. $\square$

\medskip

\noindent We note that Theorem 11.3 corresponds to the special case of Theorem 11.6
where $T(x,\sigma)=\sigma T_1x$. Now let $Q$ be an operator of the form (11.3) and suppose
that $Q$, $\Omega$ and $\varphi$ satisfy the hypotheses of Theorem 11.4. In order to apply Theorem
11.6 to the Dirichlet problem, $Qu=0$ in $\Omega$, $u=\varphi$ on $\partial\Omega$, we imbed this problem
in a family of problems,

\[
Q_{\sigma}u=a^{ij}(x,u,Du;\sigma)D_{ij}u+b(x,u,Du;\sigma)=0\ \text{in }\Omega,
\]
\[
u=\sigma\varphi\ \text{on }\partial\Omega,\qquad 0\leq \sigma\leq 1,
\]

\noindent such that:

\begin{align*}
\text{(i)}\ & Q_1=Q,\ b(x,z,p;0)=0;\\
\text{(ii)}\ & \text{the operators }Q_{\sigma}\text{ are elliptic in }\Omega\text{ for all }\sigma\in[0,1];\\
\text{(iii)}\ & \text{the coefficients }a^{ij},\, b\in C^0\!\bigl(C^{\alpha}(\overline{\Omega}\times \R\times \R^n);[0,1]\bigr),\ \text{that is, }a^{ij},\, b\in\\
& C^{\alpha}(\overline{\Omega}\times \R\times \R^n)\text{ for each }\sigma\in[0,1]\text{ and considered as mappings from }[0,1]\\
& \text{into }C^{\alpha}(\overline{\Omega}\times \R\times \R^n),\ \text{the functions }a^{ij},\, b\text{ are continuous.}
\end{align*}

\noindent For all $v\in C^{1,\beta}(\overline{\Omega})$, $\sigma\in[0,1]$ the operator $T$ is defined by letting $u=T(v,\sigma)$ be the
unique solution in $C^{2,\alpha\beta}(\overline{\Omega})$ of the linear Dirichlet problem,

\[
a^{ij}(x,v,Dv;\sigma)D_{ij}u+b(x,v,Dv;\sigma)=0\ \text{in }\Omega,\qquad u=\sigma\varphi\ \text{on }\partial\Omega.
\]

\noindent From condition (i) above we see that the solvability of the Dirichlet problem,
$Qu=0$ in $\Omega$, $u=\varphi$ on $\partial\Omega$, in the space $C^{2,\alpha}(\overline{\Omega})$ is equivalent to the solvability of
the equation $u=T(u,1)$ in the Banach space $C^{1,\beta}(\overline{\Omega})$, and that $T(v,0)=0$ for all
$v\in C^{1,\beta}(\overline{\Omega})$. The continuity and compactness of the mapping $T$ are assured by
conditions (ii) and (iii); the details of this argument are similar to the proof of
Theorem 11.4 and are left to the reader. Hence we can conclude from Theorem 11.6
the following generalization of Theorem 11.4.

\bigskip

\noindent \textbf{Theorem 11.8.} \ \emph{Let $\Omega$ be a bounded domain in $\R^n$ with boundary $\partial\Omega\in C^{2,\alpha}$ and let
$\varphi\in C^{2,\alpha}(\overline{\Omega})$. Let $\{Q_{\sigma},\ 0\leq \sigma\leq 1\}$ be a family of operators satisfying conditions
(i), (ii), (iii) above and suppose that for some $\beta>0$ there exists a constant $M$, inde-
pendent of $u$ and $\sigma$, such that every $C^{2,\alpha}(\overline{\Omega})$ solution of the Dirichlet problems
$Q_{\sigma}u=0$ in $\Omega$, $u=\sigma\varphi$ on $\partial\Omega$, $0\leq \sigma\leq 1$, satisfies}

\[
\|u\|_{C^{1,\beta}(\overline{\Omega})}<M.
\]

\noindent \emph{Then the Dirichlet problem, $Qu=0$ in $\Omega$, $u=\varphi$ on $\partial\Omega$, is solvable in $C^{2,\alpha}(\overline{\Omega})$.}


% PDF page 300
\source{gilbarg-trudinger:page-0300}

\section*{11.5. Variational Problems}

We note here that by means of the theory of topological degree (see [LS]), the hypotheses of Theorems 11.6 and 11.8 can be weakened slightly. However, since the improvement so obtained is not relevant to the particular applications in this book, we have chosen to avoid the theory of degree altogether.

In this section we consider variational problems and, in particular, their relation-
ship with elliptic partial differential equations. Let $\Omega$ be a bounded domain in $\mathbf{R}^n$
and $F$ a given function in $C^1(\Omega \times \mathbf{R} \times \mathbf{R}^n)$. We consider the functional $I$ defined
on $C^{0,1}(\overline{\Omega})$ by

\begin{equation}
\tag{11.18}
I(u)=\int_{\Omega} F(x, u, Du)\,dx.
\end{equation}

Note that, since $u \in C^{0,1}(\overline{\Omega})$, the gradient $Du$ exists almost everywhere and is
measurable and bounded; (see Section 7.3). Now let $\varphi$ be a given $C^{0,1}(\overline{\Omega})$ function
and consider $I(u)$ for all functions $u$ in the set

\[
\mathcal{C}=\{u \in C^{0,1}(\overline{\Omega}) \mid u=\varphi \text{ on } \partial\Omega\}.
\]

The problem we now consider is the following:

$\mathscr{P}$: Find $u \in \mathcal{C}$ such that $I(u) \leq I(v)$ for all $v \in \mathcal{C}$.

Let us suppose that $u$ is a solution of $\mathscr{P}$ and let $\eta$ belong to the space

\[
\mathcal{C}_0=\{\eta \in C^{0,1}(\overline{\Omega}) \mid \eta=0 \text{ on } \partial\Omega\};
\]

then the function $v=u+t\eta$ must belong to $\mathcal{C}$ for every $t \in \mathbf{R}$. Thus $I(u) \leq I(u+t\eta)$
for all $t \in \mathbf{R}$, or defining $\mathscr{J}(t)=I(u+t\eta)$, we have $\mathscr{J}(0) \leq \mathscr{J}(t)$ for all $t \in \mathbf{R}$, and so
$\mathscr{J}$ has a minimum at $0$, whence $\mathscr{J}'(0)=0$. By differentiation, we therefore obtain
the equation

\begin{equation}
\tag{11.19}
\int_{\Omega} \{D_{p_i}F(x, u, Du)D_i\eta + D_zF(x, u, Du)\eta\}\,dx=0
\end{equation}

for all $\eta \in \mathcal{C}_0$, that is the function $u$ is a weak solution of the \emph{Euler-Lagrange}
equation

\begin{equation}
\tag{11.20}
Qu=\operatorname{div} D_pF(x, u, Du)-D_zF(x, u, Du)=0.
\end{equation}

Moreover if $F \in C^2(\Omega \times \mathbf{R} \times \mathbf{R}^n)$ and $u \in C^2(\Omega)\cap C^{0,1}(\overline{\Omega})$, then $u$ is a solution

% PDF page 301
\source{gilbarg-trudinger:page-0301}

\pagestyle{empty}


of the classical Dirichlet problem $Qu=0$ in $\Omega$ in $\Omega$, $u=\varphi$ on $\partial\Omega$. Hence the solvability of problem $\mathcal{P}$ implies the solvability of the Dirichlet problem for equation (11.20).

\medskip

Let us call the functional $I$ regular if the integrand $F$ is strictly convex with respect to the $p$ variables. Clearly if $F\in C^2(\Omega\times\mathbb{R}\times\mathbb{R}^n)$, then the regularity of $I$ is equivalent to the ellipticity of the Euler-Lagrange operator $Q$. We suppose now that the function $u\in C^{0,1}(\Omega)$ satisfies (11.20) and $u=\varphi$ on $\partial\Omega$. Then we have

\[
\mathcal{P}(t)=\mathcal{P}(0)+t\,\mathcal{P}'(0)+\frac{t^2}{2}\,\mathcal{P}''(\zeta)
\]
\[
=\mathcal{P}(0)+\frac{t^2}{2}\,\mathcal{P}''(\zeta)
\]

for some $\zeta$ such that $|\zeta|\leq |t|$. If we now assume that the function $F$ is jointly convex in $z$ and $p$ so that the matrix

\[
\begin{bmatrix}
D_{p_i p_j}F & D_{p_i z}F \\
D_{p_j z}F & D_{zz}F
\end{bmatrix}
\]

is non-negative in $\Omega\times\mathbb{R}\times\mathbb{R}^n$, we have

\[
\mathcal{P}''(\zeta)=\int_{\Omega}\{D_{p_i p_j}F(x,u+\zeta\eta,Du+\zeta D\eta)D_i\eta\,D_j\eta
\]
\[
+2D_{p_i z}F(x,u+\zeta\eta,Du+\zeta D\eta)D_i\eta
\]
\[
+D_{zz}F(x,u+\zeta\eta,Du+\zeta D\eta)\eta^2\}\,dx\geq 0,
\]

and hence $\mathcal{P}(0)\leq \mathcal{P}(t)$ for all $t\in\mathbb{R}$. Consequently the function $u$ is a solution of the variational problem $\mathcal{P}$ and moreover if $I$ is regular, we see from Theorem 10.1 that $u$ is uniquely determined. Therefore we have proved

\medskip

\noindent\textbf{Theorem 11.9.}\quad \textit{Let $I$ be regular with $F$ jointly convex in $z$ and $p$. Then the variational problem $\mathcal{P}$ can have at most one solution. Furthermore, the solvability of $\mathcal{P}$ is equivalent to the solvability of the Dirichlet problem for the Euler-Lagrange equation, $Qu=0$, $u=\varphi$ on $\partial\Omega$, in the space $C^{0,1}(\overline{\Omega})$.}

\medskip

\noindent\textbf{Alternative Approaches}

By utilizing direct procedures in the calculus of variations, we can develop alternative approaches to the Dirichlet problem for variational operators $Q$. Direct methods which involve the enlargement of the set $\mathcal{E}$ to a subset of an appropriate space of weakly differentiable functions are treated in [LU 4] and [MY 5]. We describe briefly a further approach which has the advantages that the integrand $F$ need not be $C^2$ and that the solutions obtained are automatically in $C^{0,1}(\overline{\Omega})$. Let


% PDF page 302
\source{gilbarg-trudinger:page-0302}

us define, for $K\in\mathbb{R}$,

\[
\mathcal{C}_K=\{u\in \mathcal{C}\mid \|u\|_{C^{0,1}(\Omega)}\le K\}
\]

and consider the problem

\[
\mathscr{P}_K:\ \text{Find }u\in \mathcal{C}_K\text{ such that }I(u)\le I(v)\text{ for all }v\in \mathcal{C}_K.
\]

For $\mathscr{P}_K$ we have the following existence result.

\medskip

\noindent\textbf{Theorem 11.10.} Let $F\in C^1(\Omega\times \mathbb{R}\times \mathbb{R}^n)$ and suppose that $F$, $D_zF$, $D_{p_i}F\in C^0(\overline{\Omega}\times \mathbb{R}\times \mathbb{R}^n)$, $i=1,\ldots,n$. Then, if $F$ is convex with respect to $p$, the problem $\mathscr{P}_K$ is solvable for any $K$ such that $\mathcal{C}_K$ is non-empty.

\bigskip

\noindent\textit{Proof.} We show that the functional $I$ is lower semicontinuous on $\mathcal{C}_K$ with respect to uniform convergence in $\Omega$. Since $I$ is also bounded below on $\mathcal{C}_K$ and $\mathcal{C}_K$ is precompact in $C^0(\overline{\Omega})$, the result follows. Hence let $\{u_m\}\subset \mathcal{C}_K$ converge uniformly to a function $u\in \mathcal{C}_K$. We then have

\begin{equation*}
\begin{aligned}
&\text{(11.21)}\qquad I(u_m)-I(u)
&= \int_{\Omega}\bigl[F(x,u_m,Du_m)-F(x,u,Du)\bigr]\,dx \\
&= \int_{\Omega}\bigl[F(x,u_m,Du_m)-F(x,u,Du_m)\bigr]\,dx \\
&\quad + \int_{\Omega}\bigl[F(x,u,Du_m)-F(x,u,Du)\bigr]\,dx \\
&\ge -\sup_{\Omega\times \mathcal{C}_K}|D_zF|\int_{\Omega}|u_m-u|\,dx \\
&\quad + \int_{\Omega} D_{p_i}F(x,u,Du)D_i(u_m-u)\,dx
\end{aligned}
\end{equation*}

by the convexity of $F$ with respect to $p$. For fixed $i$, let us set $\varphi=D_{p_i}F(x,u,Du)$ and suppose first that $\varphi\in C^1_0(\Omega)$. Integrating by parts, we then have

\[
\int_{\Omega}\varphi D_i(u_m-u)\,dx=-\int_{\Omega}(u_m-u)D_i\varphi\,dx\to 0\quad \text{as }m\to\infty.
\]

If $\varphi\notin C^1_0(\Omega)$, then since $\varphi\in L^\infty(\Omega)$ there exists, for $\epsilon>0$, a function $\varphi_\epsilon\in C^1_0(\Omega)$ such that

\[
\int_{\Omega}|\varphi_\epsilon-\varphi|\,dx<\frac{\epsilon}{2K}.
\]

% PDF page 303
\source{gilbarg-trudinger:page-0303}


Then we have
\[
\left|\int_{\Omega} \varphi\, \Dx{i}(u_m-u)\,dx\right|
\le
\left|\int_{\Omega} \varphi_{\epsilon}\,\Dx{i}(u_m-u)\,dx\right|
+ \int_{\Omega} |\varphi_{\epsilon}-\varphi|\,|\Dx{i}(u_m-u)|\,dx.
\]

But
\[
\int_{\Omega} \varphi_{\epsilon}\,\Dx{i}(u_m-u)\,dx \to 0
\quad \text{as } m \to \infty,
\]
and as $u_m, u \in \CalK$, $|\Dx{i}(u_m-u)|\le 2K$, so that
\[
\int_{\Omega} |\varphi_{\epsilon}-\varphi|\,|\Dx{i}(u_m-u)|\,dx < \epsilon.
\]

Hence
\[
\limsup_{m\to\infty}\left|\int_{\Omega} \varphi\,\Dx{i}(u_m-u)\,dx\right|\le \epsilon
\]

and, since $\epsilon$ can be chosen arbitrarily, we conclude from (11.21) that
\[
\liminf_{m\to\infty} I(u_m)\ge I(u),
\]
that is, $I$ is lower semicontinuous on $\CalK$ with respect to uniform convergence.
\hfill $\square$

Let us call a solution of problem $\CalP_K$, a $K$-quasisolution of problem $\CalP$.
If there exists some space in which the family of all $K$-quasisolutions for
$K\in\mathbb{R}$ is relatively compact, we can obtain a generalized solution of
problem $\CalP$ as the limit of a sequence of quasisolutions $\{u_m\}$ corresponding
to constants $\{K_m\}$, $K_m\to\infty$. The following theorem shows that the
solvability of problem $\CalP$, as formulated, is a consequence of an apriori
bound in $C^{0,1}(\overline{\Omega})$ for quasisolutions.

\textbf{Theorem 11.11.} \emph{Let $u$ be a $K$-quasisolution of problem $\CalP$ such that}
\begin{equation}
\|u\|_{C^{0,1}(\overline{\Omega})} < K.
\tag{11.22}
\end{equation}
\emph{Then, if $F \in C^{1}(\Omega\times\mathbb{R}\times\mathbb{R}^n)$ is jointly convex with respect to $z$ and $p$, the function $u$ is also a solution of problem $\CalP$.}

\emph{Proof.} \ Let $v \in \CalE$. Then by (11.22) we have
\[
w=u+\epsilon(v-u)\in \CalK
\]

% PDF page 304
\source{gilbarg-trudinger:page-0304}

for some $\varepsilon>0$. Since $u$ solves $\mathscr{P}_k$, we have

\[
\int_{\Omega} F(x,u,Du)\,dx \leq \int_{\Omega} F(x,w,Dw)\,dx,
\]

but, as $w=(1-\varepsilon)u+\varepsilon v$ and $F$ is convex in $(z,p)$,

\[
\int_{\Omega} F(x,w,Dw)\,dx \leq (1-\varepsilon)\int_{\Omega} F(x,u,Du)\,dx+\varepsilon\int_{\Omega} F(x,v,Dv)\,dx.
\]

Consequently

\[
\int_{\Omega} F(x,u,Du)\,dx \leq \int_{\Omega} F(x,v,Dv)\,dx.\quad \square
\]

The combination of Theorems 11.10 and 11.11 can be viewed as analogous to Theorems 11.4 and 11.8. It is practicable to carry out the required estimation of quasisolutions in three steps corresponding to steps (i), (ii) and (iii) of the existence procedure described in Section 11.3. Namely:

\[
\text{(i)'}\ \text{Estimation of }\sup_{\Omega}|u|;
\]

\[
\text{(ii)'}\ \text{Using (i)', estimation of}
\]

\[
l'(u)=\sup_{x\in\Omega,\,y\in\partial\Omega}\frac{|u(x)-u(y)|}{|x-y|};
\]

\[
\text{(iii)'}\ \text{Using (ii)', estimation of}
\]

\[
l(u)=\sup_{x,y\in\Omega}\frac{|u(x)-u(y)|}{|x-y|}.
\]

It turns out that many of our estimates in Chapters 10, 15 and 16, (in particular the comparison principle, Theorem 10.7), can be adapted to hold for quasisolutions of variational problems, thereby facilitating the above steps. Furthermore, under appropriate hypotheses on $Q$ and $\partial\Omega$, we can obtain by regularity considerations classical solutions of the Dirichlet problem $Qu=0$, $u=\varphi$ on $\partial\Omega$.

We note here also that the methods described above can be extended to the class of divergence structure operators using the theory of monotone operators and also to encompass obstacle problems. In these situations the solvability of problem $\mathscr{P}_k$ is generalized to that of a variational inequality. The reader is referred to the works [BW 3], [HS], [LL], [LST], [PA], [WL], [KST] for further information.

% PDF page 305
\source{gilbarg-trudinger:page-0305}

\begin{quote}
\textbf{Notes}
\end{quote}

The Schauder fixed point theorem, Theorem 11.1, was established in [SC 1] and applied by Schauder to nonlinear equations in [SC 3]. Theorems 11.3 and 11.6 are special cases of the Leray-Schauder theorem [LS]. Our proofs of these results follow respectively those of Schaefer [SH] and Browder [BW 2]. The application to the Dirichlet problem, Theorem 11.5, is adapted from Gilbarg [GL 2]. For Equation (11.7) Morrey ([MY 5] p. 98) proves Theorem 11.5 assuming only the bounded slope condition, without the regularity hypotheses on $\Omega$ and $\varphi$ in Theorem 11.4.

In the first edition of this work we proved the Brouwer fixed point theorem following [DS]. In recent years many elegant and simple proofs have appeared in the literature.

% PDF page 306
\source{gilbarg-trudinger:page-0306}

Chapter 12

Equations in Two Variables

The theory of quasilinear elliptic equations in two dimensions is in many respects simpler and in some respects more general than that in higher dimensions. This chapter is concerned with aspects of the theory that are specifically two-dimensional in character, although the basic results on quasilinear equations can be extended to higher dimensions by other methods. As will be seen, the special features of this theory are founded on strong apriori estimates that are valid for general \emph{linear} equations in two variables.

12.1. Quasiconformal Mappings

Various function theoretic concepts and methods play a special role in the theory of elliptic equations in two variables ; (see [CH], for example). Here we shall be concerned mainly with apriori estimates arising from the theory of quasiconformal mappings. A continuously differentiable mapping $p=p(x,y)$, $q=q(x,y)$ from a domain $\Omega$ in the $z=(x,y)$ plane to the $w=(p,q)$ plane is \emph{quasiconformal}, or $K$-\emph{quasiconformal}, in $\Omega$ if for some constant $K>0$ we have

\begin{equation}
p_x^2+p_y^2+q_x^2+q_y^2 \leq 2K(p_xq_y-p_yq_x)
\tag{12.1}
\end{equation}

for all $(x,y)\in\Omega$. Although it suffices for the present purposes that $p$ and $q$ are in $C^1(\Omega)$, the results developed in this section will apply as well to continuous $p,q$ in $\Wloc$, that is, to continuous $p,q$ that have locally square integrable weak derivatives.

When $K<1$, (12.1) is seen to imply that $p$ and $q$ are constant and we shall therefore assume $K\ge 1$. For $K=1$, the mapping $w(z)=p(z)+iq(z)$ defines an analytic function of $z$. When $K\ge 1$, the inequality (12.1) has the geometric meaning that at points of non-vanishing Jacobian the mapping between the $z$ and $w$ planes preserves orientation and takes infinitesimal circles into infinitesimal ellipses of uniformly bounded eccentricity, in which the ratio of minor to major axis is bounded below by $\alpha=K-(K^2-1)^{1/2}>0$. This remark can be verified by direct calculation.

It will also be of interest to consider the more general class of mappings $(x,y)\to$

$(p,q)$ defined by the inequality
